% SPDX-FileCopyrightText: 2016-2026 Lari Natri <lari.natri@iki.fi>
% SPDX-License-Identifier: GPL-3.0-or-later
%
% Song navigation, indentation, alternate-chord declarations and lyric helpers.
% This file is part of the 'ulsbs' package.
%
% Load: ulsbs.sty, immediately after front matter in the song-content phase.
% Requires: typography, final style, book policy, xcolor, etoolbox and xparse.
% State: owns alternate-chord color and capture declarations; inherits @ as a
% letter without an enclosing group. Expl3 syntax switches remain local to
% their original declaration spans. Color initialization executes on load;
% song lifecycle code later resets it and the chord owner installs capture.
% Verse spacing and translation state are supplied before document-time use.
% Public providecommand style defaults may be defined before loading or renewed
% afterward; styles and indentation dimensions are read when helpers are used.

\ProvidesFile{ulsbs-song-helpers.code.tex}[ULSBS song and lyric helpers]

% === Navigation and verse indentation helpers ===

% Vertical spacing uses vspaceifnotinverse from the verse owner.
\newcommand{\ulsbs@gotovspace}{.618034\baselineskip}

% These commands will display some text followed by three dots, and aligned to
% the right of the page. The idea is to use it as a signal for the reader to
% go to a different part of the song, for example the chorus. You can use it
% on the last line of a verse or between verses (in such a case it will be put
% on its own line). The only parameter is to be the first words of the lyrics
% to go to.
%
% There are four versions of the command, each taking different number of
% parameters. It is recommended to use the one-parameter version \goto
% and if needed, use it on consecutive lines to signal consecutive jumps.
% But, if for space reasons it is deemed necessary, one can use the other
% versions to squeeze more than one successive jump on one line.
\newcommand{\goto}[1]{%
  \hfill\ulsbs@gtctext{#1}\par%
}%
\newcommand{\gotoii}[2]{%
  \vspaceifnotinverse{\ulsbs@gotovspace}%
  \hfill\ulsbs@gtctext{#1} \ulsbs@gtctext{#2}\par%
}
\newcommand{\gotoiii}[3]{%
  \vspaceifnotinverse{\ulsbs@gotovspace}%
  \hfill\ulsbs@gtctext{#1} \ulsbs@gtctext{#2} \ulsbs@gtctext{#3}\par%
}
\newcommand{\gotoiv}[4]{%
  \vspaceifnotinverse{\ulsbs@gotovspace}%
  \hfill\ulsbs@gtctext{#1} \ulsbs@gtctext{#2} \ulsbs@gtctext{#3} \ulsbs@gtctext{#4}\par%
}
% Used by \goto?? commands:
\newcommand{\ulsbs@gtctext}[1]{{\normalsize\bfseries{\ulSymbolTwoFont➞} \textit{#1}\ldots}}

% This command will make a small horizontal space. Use this in the beginning
% of each line you want to indent. The optional parameter specifies the indent
% depth factor. The default is 1.
%
% To indent complete verses, use \beginverse[<n>] instead.
\newcommand\ind[1][1]{%
  \leavevmode\kern\dimexpr #1\ulVerseIndentWidth\relax%
}

% This command puts an arrow with a text "next song" above it to the end of
% the current line. It is meant to be used on the last line of the last
% verse, to suggest that one can continue playing the next song straight.
\newcommand{\tonextsongmark}[0]{
  \hfill%
  \makebox[0pt][l]{\raisebox{1ex}{%
    \tiny\textit{\ulIfLanguage{portuguese}{seguinte}{next song}}%
  }}%
  $\Longrightarrow$
}


% === Alternate chords ===

% Initialize the color variable.
\definecolor{ulsbs@currentaltchordcolor}{rgb}{0,0,0}%
% Set default color to ulAltChordColorA, will also be reset to this at each
% song's end.
\colorlet{ulsbs@currentaltchordcolor}{ulAltChordColorA}%

% Set the color for displaying alternate chords, given as a number 1..6.
% The color will be reset to 1 at the beginning of the document, and at the
% end of each song. Called from \ulsbs@altchordsid and the song lifecycle.
\newcommand{\ulsbs@setaltchordscolor}[1]{%
  \ifnum#1=1%
    \xglobal\colorlet{ulsbs@currentaltchordcolor}{ulAltChordColorA}%
  \fi%
  \ifnum#1=2%
    \xglobal\colorlet{ulsbs@currentaltchordcolor}{ulAltChordColorB}%
  \fi%
  \ifnum#1=3%
    \xglobal\colorlet{ulsbs@currentaltchordcolor}{ulAltChordColorC}%
  \fi%
  \ifnum#1=4%
    \xglobal\colorlet{ulsbs@currentaltchordcolor}{ulAltChordColorD}%
  \fi%
  \ifnum#1=5%
    \xglobal\colorlet{ulsbs@currentaltchordcolor}{ulAltChordColorE}%
  \fi%
  \ifnum#1=6%
    \xglobal\colorlet{ulsbs@currentaltchordcolor}{ulAltChordColorF}%
  \fi%
}

% [set number]{id text}: visible IDs change the alternate-set color until the
% song ends or another ID changes it. Hidden arguments must remain opaque, even
% to TeX's conditional scanner, and must not change that persistent color.
\ExplSyntaxOn
\NewDocumentCommand \ulsbs@altchordsid {O{1} m}
  {\bool_if:nT {\legacy_if_p:n {ulShowChords} && \legacy_if_p:n {ulShowAltChords}}
   {
    \ulsbs@setaltchordscolor{#1}%
    \color{ulsbs@currentaltchordcolor}%
    \tl_if_empty:nF {#2}{
      \renewcommand\capo[1]{%
        {\normalfont (capo~##1)}%
      }%
      \makebox[0pt][r]{\raisebox{3.75ex}[0pt][0pt]{%
        \tiny%
        \ulsbs@altchordidbox{\color{ulsbs@currentaltchordcolor}\normalfont\textbf{#2}:\ulsbs@vphantommaxtext}%
      }}%
    }
   }}
\ExplSyntaxOff

% Consumer override: define this providecommand default before loading or renew
% after loading. It is added to the normal chord style at each use.
\providecommand{\ulAltChordsStyle}{%
  \footnotesize\itshape%
}

\ExplSyntaxOn
% True if within \altchords parameter
\bool_new:N \l__ulsbs_chord_in_altchords_bool


% Meant to be used at the end of a lyric line to specify an alternate set
% of chords for that line to be displayed floated to the right of the line
% in question. Do NOT set chords with \[] macro inside this, as it would mess
% with the chord repeats etc. Instead use plain text for chord names. They are
% transposed only when \ulTransposeAltChords{true} is set. Not displayed on books
% without chords. Only place bar lines, plain text, \id, \ac, \notrans and \e
% within the parameter. Sharps and flats may be written as # and &, or with
% the retained \shrp and \flt macros.
%
% To show an ID text, use \id macro within \altchords, preferably as the first
% thing on the first line of this chord set. The optional parameter of \id sets
% the color 1..6 for this chord set, which will be held until the end of the
% song. Numbers 1..3 are meant to be used for different keys with same chords,
% and numbers 4..6 are for alternate chordings in the default key.
% See \ulsbs@altchordsid.
%
% Example:
%
%  \altchords{\id[1]{Peter (C)}|C | |G |C#}
%
% The chord owner installs argument capture/catcodes and the lazy display guard.
\newcommand{\altchords}{}
\ExplSyntaxOff

% === Lyric helpers ===

% Just prints a dash if measure bar lines are enabled and does nothing otherwise.
% Use this if a line would end at a new measure starting without any lyrics in that
% bar on the same line. This signifies that there indeed is a measure here. The
% measure might or might not continue on the next line. You can think of this as
% an abbreviation for "Empty" or "End". :)
\newcommand{\e}{\ifulShowMeasures--\fi}

% Simply a shortened form for \textsuperscript
\newcommand{\up}[1]{%
  \textsuperscript{#1}%
}

\ExplSyntaxOn
% Print an alternate chord for a numbered repetition.
% Usage: \ac{<repetition>}{<chord>}, for example \ac{2}{Am}.
% Modules install these declared public command implementations below.
\NewDocumentCommand{\ac}{m m}{}
\ExplSyntaxOff

% Display alternate lyrics, but without any vertical or horizontal space. Is called
% from \altlyr
\newcommand{\altlyrnospace}[2][*]{%
  {\ifulsbs@intranslation\else\small\fi\up{#1}\emph{#2}}%
}

% Display alternate lyrics. If inside a verse, show them on the same line
% (pushed to the right), otherwise include a small vertical space before.
% Usage \altlyr[<superscript>]{<lyrics>}
% If the optional argument is omitted, it will be replaced with an asterisk.
% Example: \altlyr[2]{with love}
\newcommand{\altlyr}[2][*]{%
  \vspaceifnotinverse{.618034\baselineskip}%
  \hfill%
  \altlyrnospace[#1]{#2}%
}

% Put tiny sublyrics, given as the parameter, below the main lyrics (could be
% used for a second voice). This is not very robust; a good idea might be to
% put each bar's sublyrics as it's own \sublyr command. The sublyrics are
% left-aligned with whatever comes immediately after the command, which might
% cause overlapping of \sublyr text with the previous one.
\newcommand{\sublyr}[1]{%
  \makebox[0pt][l]{\raisebox{-1ex}[0pt][1ex]{%
    \tiny #1%
  }}%
}

% Same as \sublyr, but this pushes the main lyrics to the right to make room
% for the sublyrics, instead of aligning. Use carefully.
\newcommand{\sublyrpush}[1]{%
  \makebox{\raisebox{-1ex}[0pt][1ex]{%
    \tiny #1%
  }}%
}
